% Proposed section for the ProveIt Polylogarithms manuscript.
% Parent: AMS mathematics; theorem, lemma, corollary, remark; hyperref.
% Full proofs and provenance: accompanying article.tex, sections 2--6.
% Not an automatic replacement of any existing chapter.
\providecommand{\corrclosureG}{\mathcal G}
\providecommand{\corrclosureK}{\mathcal K}
\providecommand{\corrclosureA}{\mathcal A_{\mathrm{cc}}}
\providecommand{\corrclosureB}{\mathcal B_{\mathrm{cc}}}
\providecommand{\corrclosureFP}{\operatorname{FP}}
\providecommand{\corrclosureLi}{\operatorname{Li}}
\section{Bilinear finite-part Stieltjes correlations}
\label{corrclosure:sec}

Let $\zeta(1+u,x)=u^{-1}+\sum_{j\ge0}(-1)^j\gamma_j(x)u^j/j!$,
so $\gamma_0=-\psi$. For $0<a<1$ put $b=1-a$ and define
\begin{align}
 I_{mn}(a)=\lim_{\epsilon\downarrow0}\biggl\{
 &\int_\epsilon^b\gamma_m(x)\gamma_n(a+x)\,dx
 +\int_\epsilon^a\gamma_m(b+y)\gamma_n(y)\,dy\notag\\
 &+\frac{\gamma_n(a)}{m+1}\log^{m+1}\epsilon
 +\frac{\gamma_m(b)}{n+1}\log^{n+1}\epsilon\biggr\}.
 \label{corrclosure:definition}
\end{align}
This is the finite part of a circular translation, not an ordinary
integral. The identity
$\gamma_m(x)=x^{-1}\log^m x+\gamma_m(1+x)$ proves that the limit exists:
subtract the two displayed local singular terms inside the integrals,
then add their integrated upper-endpoint terms. The resulting integrals
are absolutely convergent.

\begin{theorem}[All-order bilinear closure]
\label{corrclosure:theorem}
For $m,n\ge0$ and $N=m+n+1$,
\begin{align}
 I_{mn}(a)={}&\frac{\gamma_N(a)}{m+1}
     +\frac{\gamma_N(b)}{n+1}+e_{mn}\notag\\
 &+\sum_{j=0}^{N-2}\bigl(c_{mn,j}\gamma_j(a)
                             +d_{mn,j}\gamma_j(b)\bigr),
 \label{corrclosure:closure}
\end{align}
where the sum is empty for $N=1$. All coefficients are explicit
polynomials over $\mathbb Q$ in $\zeta(2),\ldots,\zeta(N+1)$.
There is no $\gamma_{N-1}$ term. Also $I_{mn}(a)=I_{nm}(1-a)$.
\end{theorem}
\begin{proof}
Put $w=u+v$ and
\[
 \corrclosureA(u,v)=\frac{\Gamma(1+w)\Gamma(1-v)}{\Gamma(1+u)},
 \qquad \corrclosureB(u,v)=\corrclosureA(v,u),\qquad
 \corrclosureG(u,x)=\zeta(1+u,x)-\frac1u.
\]
Classical Hurwitz Fourier orthogonality, initially in an absolutely
convergent spectral region, and Gamma reflection give the kernel
\begin{equation}
 \corrclosureK=-u\corrclosureA(u,v)\zeta(1+w,b)
                    -v\corrclosureB(u,v)\zeta(1+w,a).
 \label{corrclosure:kernel}
\end{equation}
It is $uv$ times the meromorphic continuation of the circular product
integral of $\zeta(1+u,x)$ and $\zeta(1+v,x\mathbin\oplus a)$.
For an explicit endpoint justification, that product's regular part is
\begin{align*}
 R={}&\int_0^b\{x^{-1-u}[\corrclosureG(v,a+x)-\corrclosureG(v,a)]
                   +\corrclosureG(u,1+x)\corrclosureG(v,a+x)\}\,dx\\
 &+\int_0^a\{y^{-1-v}[\corrclosureG(u,b+y)-\corrclosureG(u,b)]
                   +\corrclosureG(v,1+y)\corrclosureG(u,b+y)\}\,dy\\
 &+\corrclosureG(v,a)\frac{1-b^{-u}}u
      +\corrclosureG(u,b)\frac{1-a^{-v}}v.
\end{align*}
The quotients take their removable values $\log b$ and $\log a$.
The differences vanish linearly at the respective endpoints, so $R$
is holomorphic near zero, including after spectral differentiation.
Its coefficients are exactly $(-1)^{m+n}I_{mn}/(m!n!)$.
Using zero mean of the Hurwitz function in the initial convergence
region identifies $R$ as the mixed divided difference of
$\corrclosureK$. Hence
\[
 I_{mn}=(-1)^{m+n}m!n![u^{m+1}v^{n+1}]\corrclosureK.
\]
Since $\corrclosureA(u,-u)=\corrclosureB(u,-u)=1$, the quotient
$Q=(u\corrclosureA+v\corrclosureB)/(u+v)$ is holomorphic, and
\[
 \corrclosureK=-Q-u\corrclosureA\corrclosureG(w,b)
                    -v\corrclosureB\corrclosureG(w,a).
\]
Finally,
\[
 \log\corrclosureA=\sum_{k\ge2}\frac{\zeta(k)}k
       \{(-1)^k((u+v)^k-u^k)+v^k\}.
\]
This proves finite linear closure and the coefficient algebra. The
constant Gamma-ratio term gives the two coefficients of $\gamma_N$;
the missing degree-one Gamma-ratio term gives the absence of
$\gamma_{N-1}$. Interchanging the variables proves symmetry.
The accompanying report supplies the full Fourier derivation and the
exact finite coefficient algorithm.
\end{proof}

\begin{corollary}[First identities and an antiderivative]
\label{corrclosure:examples}
For $0<a<1$,
\begin{align*}
 I_{00}(a)&=\gamma_1(a)+\gamma_1(1-a)-\frac{\pi^2}{3},\\
 I_{10}(a)&=\tfrac12\gamma_2(a)+\gamma_2(1-a)
        +\zeta(2)(\gamma_0(a)+\gamma_0(1-a))-\zeta(3),\\
 \int I_{00}(a)\,da
 &=\tfrac12\bigl(\zeta''(0,a)-\zeta''(0,1-a)\bigr)
                         -\frac{\pi^2}3a+\text{constant}.
\end{align*}
\end{corollary}
\begin{proof}
Expand the kernel through total degrees two and three. For the
primitive, use
$\partial_a\zeta^{(j+1)}(0,a)=(-1)^{j+1}(j+1)\gamma_j(a)$.
\end{proof}

\begin{remark}[Evidence and normalization]
The accompanying delivery proves the coincident-endpoint analogue and
its collision counterterm, an ordinary log-Gamma autocorrelation,
arithmetic traces, and higher derivative and primitive formulas.
Its exact table covers $m+n\le8$, but the theorem above holds at all
orders by proof. Numerical checks are diagnostics, not interval
certificates, and no period independence or literature priority is
asserted. The classical starting identities are documented in
\href{https://dlmf.nist.gov/25.11}{DLMF Section 25.11}.
\end{remark}
